% Pursuing Stacks (cotes 134-2 to 134-8): archive cote 134-5, site page 76 (PDF page 77), typescript page 263.
% typescript. Draft compared: ps4.tex:3857-3918.
% AI-transcribed by Claude Opus 5.5 (claude-opus-5-5, adaptive thinking, effort high), Message Batches API, batch ps2 (msgbatch_01NVwr1Yuo8qJf5rykeGBXPK), 2026-11-12; not yet checked by eye.
\scanpage{134-5/76}
\margin{à la ligne}
\mspage{263}

This relationship has been "in the air" \xout{xxxx} since section 50 (p. 109-%
110),  and I kind of turned around it consistently up to section 60,
without really getting to the core.  One reason for this "turning around"
has been,  I guess,  that I let myself be distracted,  not to say hypno-%
tized,  by the"canonical" $\hand{\pi}_0$ functor on a category (which makes really
good sense only when M is"\unsure{tdally|totally} 0-connected"\xout{though} as a category, \inserted{a condition which} \deleted{which}
\over{w}{\unsure{sh}}ould mean \inserted[typed]{more or less} I suppose that the 0-connected objects of M,  defined in terms
of the mere category structure of M,  define a $\underline{W}_{\unsure{\bullet|0}}$-asphericity structure
on M )\unsure{,|.}  Even after realizing (in section 59)  that one should generalize
the \deleted{definition} \inserted{description} of a contractibility structure in terms of \xout{axxx} a "con-%
nected components functor" $\hand{\pi}_0$ ,  to the case of a functor $\hand{\pi}_0$ : M $\hand{\rightarrow}$(Sets)
given beforehand,  and satisfying suitable restrictions (which I did not
try to elucidate),  I still was under the impression\xout{s} that the contracti-%
bility structures one could get this \xout{xx} way must be of an extremely speci-%
al nature.  In order to become aware of the fact that this is by no means
so,  namely that any contractibility structure could be obtained from a
suitable functor $\hand{\pi}_0$ ,  it \xout{sh}would have been necessary to notice that such
a structure $M_c$  defines in a natural way a functor $\hand{\pi}_0$ .  There was \xout{indeed}
indeed the realization that $M_c$ should \inserted[typed]{allow to} define a notion of weak equivalence
(cf. page 136),  but it was'nt clearly realized by then that at the same
time as a notion of weak equivalence W in M,  we should also get a canoni%
cal functor
\[ M \longrightarrow W^{-1}M \longrightarrow \text{(Hot)} \quad , \]
namely something a lot more precise still than a functor with values in
(Sets)  merely !  But rather than push ahead in this direction, I then
decided (p.138)  that it would be "unreasonable" \xout{xxx} to go on \inserted[typed]{still longer} pushing \unsure{off|of}
investigation of test functors with \unsure{values|valuas} in (Cat),  following \over{an}{the} app-%
roach which had \inserted{\deleted{\illegible{1}}} been on my mind for quite a while by then,  and \deleted{written} \inserted{finally}
\deleted{down} \inserted{sketched} (with the promise of a corresponding generalization of the former
"key result" on test\unsure{n}functors)  in section 47.  Retrospectively,  the
whole "grinding" part III of these notes  now looks as a rather heavy
and \inserted[typed]{long-winded} \xout{awkwardx}digression,prompted by this approach\xout{,}  \xout{which} to still a parti-%
cular case of test functors (namely with values in (Cat),  and more strin-%
gently still \inserted{factoring through} \deleted{in} (Cat)$_{as}$ ).  \xout{which}
% NOTE: a handwritten line under the struck "which" and a double vertical stroke before "This" mark a new paragraph here
This particular case has been \xout{by} \inserted{of} no
\xout{means} use in the present part IV of the notes,  developping the really
relevant notions in terms of asphericity structures.  \uline{Technically spea\xout{k}-}%
\uline{king},  \uline{it now appears that most of the reflections of part III}  \uline{are super}%
\uline{seded by part IV}  -  the main exception being the development of the

\margin{76}
